\subsection{Weak gravitational lensing}
\label{sec:lensing}

\subsubsection{From classical scattering to light deflection}

A useful starting point is the gravitational scattering of a test particle by a point mass. This construction follows the progression in Yoo's weak-lensing lecture notes~\cite{yoo_weak_2026}: first determine the bending of a trajectory, then use the observer--lens--source geometry to relate that bending to an image position. The extension to neighboring rays will describe the distortion of a resolved galaxy.

Place a mass $M$ at the origin and let a particle approach with speed $v$ along the $z$ direction. Its unperturbed path passes the mass at transverse distance $b$, the \emph{impact parameter}. If the deflection is small, the acceleration can be evaluated along this straight path, with $z=vt$. The transverse component is
\begin{equation}
    a_{\perp}(t)=-\frac{GMb}{\left(b^2+v^2t^2\right)^{3/2}},
    \label{eq:classical-acceleration}
\end{equation}
where the negative sign indicates acceleration toward the mass. Integrating over the encounter gives
\begin{equation}
    \Delta v_{\perp}
    =-GMb\int_{-\infty}^{\infty}\frac{dt}{\left(b^2+v^2t^2\right)^{3/2}}
    =-\frac{2GM}{bv}.
    \label{eq:classical-impulse}
\end{equation}
The small bending angle is consequently $\widehat\alpha_{\mathrm{N}}\simeq |\Delta v_{\perp}|/v=2GM/(bv^2)$. This approximation requires $GM/(bv^2)\ll1$, so the trajectory used in the integral differs only slightly from the actual path~\cite{yoo_weak_2026,bartelmann_weak_2001}.

Formally setting $v=c$ gives a Newtonian estimate for light bending, but it does not derive the behavior of a null geodesic. General relativity predicts twice this angle for a ray passing a static point mass in the weak-field limit:
\begin{equation}
    \widehat\alpha=\frac{4GM}{bc^2}.
    \label{eq:point-mass-deflection}
\end{equation}
The classical analogy captures the dependence on mass and impact parameter; the additional factor of two is a relativistic result. The weak-field condition here is $GM/(bc^2)\ll1$. As the lens geometry will show, a small bending angle can still produce multiple images when the alignment is favorable~\cite{yoo_weak_2026,bartelmann_weak_2001}.

\subsubsection{Lens geometry and the Einstein radius}

For a localized lens, the \emph{thin-lens approximation} replaces the continuous bending along the ray with a change of direction at one plane. A galaxy cluster can be treated this way when its depth is small compared with the distances separating observer, lens, and source. Let $D_l$, $D_s$, and $D_{ls}$ be the angular-diameter distances from observer to lens, observer to source, and lens to source. These distances convert an angle into a transverse physical separation at the relevant plane. In an expanding universe, $D_{ls}$ is not generally $D_s-D_l$~\cite{bartelmann_weak_2001,schneider_weak_2006}.

Measure angles relative to the direction of the lens center. The observed image lies at $\boldsymbol{\theta}$, while $\boldsymbol{\beta}$ is the angular position the source would have without deflection. The ray crosses the lens plane at physical position $\boldsymbol{\xi}=D_l\boldsymbol{\theta}$. In the small-angle geometry, extending the observed ray without bending would place it at $D_s\boldsymbol{\theta}$ on the source plane. The bend changes that position by $D_{ls}\widehat{\boldsymbol{\alpha}}$, giving
\begin{equation}
    D_s\boldsymbol{\beta}
    =D_s\boldsymbol{\theta}-D_{ls}\widehat{\boldsymbol{\alpha}}(D_l\boldsymbol{\theta}).
    \label{eq:lens-geometry}
\end{equation}
For a point mass, the conventional deflection vector $\widehat{\boldsymbol{\alpha}}$ points in the direction of $\boldsymbol{\xi}$ and is subtracted in this equation. It is therefore opposite to the signed transverse velocity change in Eq.~\eqref{eq:classical-impulse}; the physical bending is toward the mass. Dividing by $D_s$ defines the scaled deflection $\boldsymbol{\alpha}$ and yields the \emph{lens equation}~\cite{yoo_weak_2026,schneider_weak_2006}:
\begin{equation}
    \boldsymbol{\beta}=\boldsymbol{\theta}-\boldsymbol{\alpha}(\boldsymbol{\theta}),
    \qquad
    \boldsymbol{\alpha}(\boldsymbol{\theta})
    =\frac{D_{ls}}{D_s}\widehat{\boldsymbol{\alpha}}(D_l\boldsymbol{\theta}).
    \label{eq:scaled-deflection}
\end{equation}

For a point mass, $b=D_l|\boldsymbol{\theta}|$. Substituting Eq.~\eqref{eq:point-mass-deflection} into the lens equation gives
\begin{equation}
    \boldsymbol{\beta}=\boldsymbol{\theta}
    -\theta_E^2\frac{\boldsymbol{\theta}}{|\boldsymbol{\theta}|^2},
    \qquad
    \theta_E^2=\frac{4GM}{c^2}\frac{D_{ls}}{D_lD_s}.
    \label{eq:einstein-radius}
\end{equation}
If the source is exactly behind the lens, $\boldsymbol{\beta}=\boldsymbol{0}$, every image direction with $|\boldsymbol{\theta}|=\theta_E$ satisfies the equation. An ideal point source then produces an Einstein ring, and $\theta_E$ is its angular radius. For an offset source, the two images lie along the source--lens axis. Using signed coordinates on that axis reduces the equation to $\beta=\theta-\theta_E^2/\theta$, with solutions
\begin{equation}
    \theta_{\pm}=\frac{\beta\pm\sqrt{\beta^2+4\theta_E^2}}{2}.
    \label{eq:point-lens-images}
\end{equation}
The opposite signs place the images on opposite sides of the lens. This example separates weak gravitational fields from weak image distortions: the field can be weak along the ray while the image mapping produces strong-lensing features near an Einstein ring~\cite{yoo_weak_2026,bartelmann_weak_2001}.

\subsubsection{An extended lens and its projected potential}

A cluster has an extended mass distribution rather than a single point mass. In the weak-field approximation, the potential and leading-order deflections of small mass elements add linearly. Summing the point-mass contributions therefore extends the scattering calculation to the whole lens. Let $\Phi(\boldsymbol{\xi},z)$ be its Newtonian gravitational potential, with $z$ now a physical line-of-sight coordinate through the lens. The deflection generalizes to the integrated transverse potential gradient,
\begin{equation}
    \widehat{\boldsymbol{\alpha}}(\boldsymbol{\xi})
    =\frac{2}{c^2}\int dz\,\boldsymbol{\nabla}_{\xi}\Phi(\boldsymbol{\xi},z).
    \label{eq:potential-deflection}
\end{equation}
The integral is evaluated along the approximately undeflected path through the lens. For $\Phi=-GM/(|\boldsymbol{\xi}|^2+z^2)^{1/2}$, its magnitude is $4GM/(c^2|\boldsymbol{\xi}|)$, recovering the point-mass result. Thus the same gradient construction applies to a continuous mass distribution~\cite{yoo_weak_2026,bartelmann_weak_2001}.

It is convenient to absorb the distance factors into a dimensionless projected lensing potential $\psi$. Since only potential differences affect deflection, choose a reference direction $\boldsymbol{\theta}_0$ away from any mass singularity and define
\begin{equation}
    \psi(\boldsymbol{\theta})=
    \frac{2D_{ls}}{c^2D_lD_s}\int dz\,
    \left[\Phi(D_l\boldsymbol{\theta},z)-\Phi(D_l\boldsymbol{\theta}_0,z)\right].
    \label{eq:projected-potential}
\end{equation}
The reference term is independent of $\boldsymbol{\theta}$ and fixes an otherwise arbitrary additive constant. It also avoids the divergent constant in the line integral of an isolated point-mass potential. The chain rule gives $\boldsymbol{\nabla}_{\theta}=D_l\boldsymbol{\nabla}_{\xi}$ when acting on the first term, so Eqs.~\eqref{eq:scaled-deflection} and~\eqref{eq:potential-deflection} imply
\begin{equation}
    \boldsymbol{\alpha}=\boldsymbol{\nabla}_{\theta}\psi,
    \qquad
    \boldsymbol{\beta}=\boldsymbol{\theta}-\boldsymbol{\nabla}_{\theta}\psi.
    \label{eq:lens-equation}
\end{equation}
Here $\Phi$ has dimensions of velocity squared, whereas $\psi$ is dimensionless and differentiated with respect to angular coordinates. These distinctions are needed when relating image distortions to a physical mass density~\cite{yoo_weak_2026,schneider_weak_2006}.

Let $\rho(\boldsymbol{\xi},z)$ be the physical density of the localized lens and define its projected surface density by $\Sigma(\boldsymbol{\xi})=\int \rho(\boldsymbol{\xi},z)\,dz$. The three-dimensional Poisson equation is $\nabla^2\Phi=4\pi G\rho$. Writing $\nabla^2=\nabla_{\xi}^2+\partial_z^2$ and integrating along the line of sight gives
\begin{equation}
    \int dz\,\nabla_{\xi}^2\Phi
    =4\pi G\Sigma
    -\left[\partial_z\Phi\right]_{-\infty}^{+\infty}
    =4\pi G\Sigma,
    \label{eq:projected-poisson}
\end{equation}
where the boundary term vanishes for an isolated lens whose gravitational field decays far from the mass. Taking the angular Laplacian of Eq.~\eqref{eq:projected-potential} then yields
\begin{equation}
    \nabla_{\theta}^2\psi
    =\frac{8\pi G}{c^2}\frac{D_lD_{ls}}{D_s}\Sigma(D_l\boldsymbol{\theta})
    =2\kappa(\boldsymbol{\theta}),
    \label{eq:lensing-poisson}
\end{equation}
where the \emph{convergence} $\kappa$ and the critical surface density are defined by
\begin{equation}
    \kappa(\boldsymbol{\theta})=\frac{\Sigma(D_l\boldsymbol{\theta})}{\Sigma_{\mathrm{crit}}},
    \qquad
    \Sigma_{\mathrm{crit}}=\frac{c^2}{4\pi G}\frac{D_s}{D_lD_{ls}}.
    \label{eq:convergence}
\end{equation}
Convergence is therefore a dimensionless measure of projected mass, scaled by the observer--lens--source geometry. For a fixed cluster, different source distances give different $\Sigma_{\mathrm{crit}}$ and hence different lensing strengths. The surface density here is mass per physical area; no comoving-area convention is being used~\cite{yoo_weak_2026,bartelmann_weak_2001}.

\subsubsection{Convergence, shear, and magnification}

A constant deflection shifts a source's apparent position. A resolved galaxy changes shape because neighboring rays are deflected by different amounts. Across an image small enough that the deflection gradient is nearly constant, expand the lens equation as $\delta\boldsymbol{\beta}=A\,\delta\boldsymbol{\theta}$. Differentiating Eq.~\eqref{eq:lens-equation} gives
\begin{equation}
    A_{ij}=\frac{\partial\beta_i}{\partial\theta_j}
    =\delta_{ij}-\psi_{,ij}
    =\begin{pmatrix}
        1-\kappa-\gamma_1 & -\gamma_2\\
        -\gamma_2 & 1-\kappa+\gamma_1
    \end{pmatrix},
    \label{eq:jacobian}
\end{equation}
where $\delta_{ij}$ is the identity matrix and commas denote angular derivatives. The trace and traceless parts of the potential Hessian define
\begin{equation}
    \kappa=\frac{\psi_{,11}+\psi_{,22}}{2},
    \qquad
    \gamma_1=\frac{\psi_{,11}-\psi_{,22}}{2},
    \qquad
    \gamma_2=\psi_{,12}.
    \label{eq:potential-derivatives}
\end{equation}
The first definition agrees with Eq.~\eqref{eq:lensing-poisson}. The equality of mixed derivatives makes $A$ symmetric: this single-plane potential mapping has no antisymmetric image-rotation term~\cite{yoo_weak_2026,schneider_weak_2006}.

The matrix $A$ maps image offsets to source offsets, so the stretching of the observed image is described by $A^{-1}$. With no shear, $A=(1-\kappa)I$; positive convergence with $\kappa<1$ enlarges the image equally in both directions. Shear produces unequal stretching. The component $\gamma_1$ distinguishes the two coordinate axes, while $\gamma_2$ describes the same type of distortion along axes rotated by \ang{45}. If $\gamma=\gamma_1+i\gamma_2$, the two eigenvalues of $A$ are
\begin{equation}
    \lambda_{\pm}=1-\kappa\pm|\gamma|.
    \label{eq:lensing-eigenvalues}
\end{equation}
A small circular source consequently appears as an ellipse with principal stretches $1/|\lambda_{\pm}|$, provided the local linear mapping is nonsingular. The determinant controls the change in angular area. Since gravitational lensing conserves surface brightness, the corresponding flux magnification is~\cite{yoo_weak_2026,schneider_weak_2006}
\begin{equation}
    \mu=\frac{1}{|\det A|}
    =\frac{1}{\left|(1-\kappa)^2-|\gamma|^2\right|}.
    \label{eq:magnification}
\end{equation}

Galaxy shapes measure the relative stretching of the two axes, rather than an overall isotropic change in size. Factoring $(1-\kappa)$ out of Eq.~\eqref{eq:jacobian} shows that these relative distortions depend on the \emph{reduced shear},
\begin{equation}
    g=g_1+ig_2=\frac{\gamma_1+i\gamma_2}{1-\kappa}.
    \label{eq:reduced-shear}
\end{equation}
The weak-lensing limit used here is $|\kappa|,|\gamma|\ll1$, in which $g\simeq\gamma$ and image distortions are small. This is more restrictive than requiring a weak gravitational field along a light ray. Cluster cores can contain strong-lensing features even when the surrounding field is suitable for weak-lensing analysis~\cite{bartelmann_weak_2001,schneider_weak_2006}.

\subsubsection{Ellipticity and shape noise}

To relate this distortion to galaxy shapes, consider an ideal image with elliptical isophotes (contours of equal surface brightness). Define its complex ellipticity by
\begin{equation}
    \epsilon=\epsilon_1+i\epsilon_2
    =\frac{a-b}{a+b}\exp(2i\varphi),
    \label{eq:ellipticity}
\end{equation}
where $a$ and $b$ are the major and minor semiaxes and $\varphi$ is the major-axis angle. The factor of two expresses a physical symmetry: rotating an ellipse by \ang{180} leaves its shape unchanged. Ellipticity and shear are consequently called spin-2 quantities. For an active rotation of the image through an angle $\alpha$ in fixed coordinates, $\epsilon$ transforms as $\epsilon\exp(2i\alpha)$~\cite{schneider_weak_2006}.

Several shape conventions appear in lensing software. In particular, the distortion $\chi=[(a^2-b^2)/(a^2+b^2)]\exp(2i\varphi)$ satisfies $\chi=2\epsilon/(1+|\epsilon|^2)$ and has a different shear response. Yoo's notes use $\epsilon$ for this quadratic-axis convention; here it is denoted $\chi$ to keep it distinct from Eq.~\eqref{eq:ellipticity}. The convention must therefore be specified before comparing outputs from different estimators. Here $\epsilon$ denotes the ideal ellipticity in Eq.~\eqref{eq:ellipticity}, before instrumental blurring, and $\boldsymbol{e}$ will denote a measured two-component shape statistic that still requires calibration~\cite{yoo_weak_2026,schneider_weak_2006}.

For a circular source in the weak-lensing regime, the two semiaxes are proportional to the reciprocal eigenvalues in Eq.~\eqref{eq:lensing-eigenvalues}. Their ratio gives $(a-b)/(a+b)=|g|$, and the orientation of the ellipse gives $\epsilon=g$. An intrinsically elliptical source already has a preferred shape before lensing. If $\epsilon^{\mathrm{s}}$ is that unlensed ellipticity, the ideal transformation for $|g|<1$ is
\begin{equation}
    \epsilon=\frac{\epsilon^{\mathrm{s}}+g}{1+g^*\epsilon^{\mathrm{s}}},
    \label{eq:ellipticity-transform}
\end{equation}
with $g^*$ the complex conjugate. Averaging over an isotropic distribution of intrinsic orientations gives $\langle\epsilon\rangle=g$. The intrinsic ellipticity of any particular galaxy is unknown, however, so its observed elongation cannot be assigned entirely to the lens. The resulting scatter is \emph{shape noise}. For $N$ independent galaxies sharing approximately the same shear, an ideal unweighted mean has per-component uncertainty $\sigma_{\epsilon,i}/\sqrt{N}$. Measurement noise, selection, and intrinsic alignments (correlations in unlensed galaxy orientations) complicate this ideal result~\cite{schneider_weak_2006,mandelbaum_weak_2018}.
